% Short subsection to include in Chapter 5 describing the I2S+filter test
\subsection{Verifikasi dengan Filter}

Pengujian tambahan dilakukan dengan menyalurkan nada sinus 1 kHz melalui jalur I2S dan
mengaplikasikan filter digital yang diimplementasikan pada lapisan perangkat lunak. Hasil
verifikasi menunjukkan bahwa filter memodifikasi konten spektral sesuai rancangan tanpa
mengganggu timing I2S yang telah diverifikasi sebelumnya (lihat Gambar \ref{fig:ch5-1khz} dan
Lampiran~\ref{app:arbitrary-fs} untuk detail desain).

\begin{figure}[htbp]
    \centering
    \IfFileExists{thesi2s/tex/figures/1khz_filtered.jpeg}{%
        \includegraphics[width=0.8\textwidth]{thesi2s/tex/figures/1khz_filtered.jpeg}%
    }{}
    \caption{Contoh sinyal 1 kHz sebelum dan setelah filter (ilustrasi).}
    \label{fig:ch5-1khz-filtered}
\end{figure}

Referensi ke Lampiran~\ref{app:arbitrary-fs} menyediakan diff kode dan penjelasan desain
untuk varian I2S yang mendukung frekuensi sampling arbitrary sebagai bonus penelitian.
